\documentclass[11pt,a4paper]{article}

\usepackage{polyglossia}
\setmainlanguage{english}

\usepackage{fsorg-manual}

\graphicspath{{figures/en/}}

\hypersetup{
  pdftitle={FS Organizer · User manual},
  pdflang={en}}

\fsorgrunningheadis{FS Organizer · User manual}

\begin{document}

\frenchspacing

\fsorgtitle
  {User manual}
  {Version \input{../VERSION.txt}}
  {Distributed under the GNU General Public License version 2.}
  {The typeface is Archivo, by Omnibus-Type, under the SIL Open Font License 1.1.}

\thispagestyle{empty}
\tableofcontents

\section{What this program does}

FS Organizer enables and disables Microsoft Flight Simulator addons without
moving their files.

You keep your addons in a folder outside the simulator, which the program calls
a \ui{library}. Its subfolders are \ui{categories}. Enabling an addon creates a
link to it in a folder the simulator reads, called a \ui{destination}.
Disabling it removes the link. The files stay in the library the whole time, so
enabling or disabling even a large addon takes no time at all.

\subsection{What the program does not store}

The program never writes down which addons are enabled. Every time it scans, it
looks at the destinations on disk and reports what is there. If you change
folders by hand, or another program's installer does, click \ui{Refresh} and the
list matches what the simulator sees.

\subsection{Profile}

A \ui{profile} is one simulator installation with its destinations and
libraries. If you have Flight Simulator 2020 and 2024 on the same computer, you
have two profiles. The program only changes the profile in use, which you
choose in \ui{Options}.

\subsection{Link type}

In \ui{Options}, on the \ui{Links} tab, you choose how links are made.

\begin{description}
  \item[Directory junction] Works across local drives without administrator
    rights. This is the default, and the one to use unless your library is on
    a network share.
  \item[Symbolic link] Only needed for a library on a network share. Windows
    only allows it with Developer Mode turned on or with administrator rights,
    and the program tells you when neither is available.
\end{description}

Changing the type affects new links only. Existing links stay as they are.

\section{Presets}

A preset is a named set of enabled addons. The same preset can be applied in
three ways, and each gives a different result, so read the plan before you
apply.

\subsection{What a preset holds}

To make a preset, enable the addons you want and click \ui{New from enabled
addons}.

The \ui{Content} column shows how many addons the preset names and how many
categories they are in. The \ui{Updated} column shows when it was last saved.

A preset can also control the simulator's startup entries
(section~\ref{sec:presets-startup}).

\subsection{A preset that matches your setup}

The \ui{If applied} column shows how many addons would change if you applied
the preset now, in the mode chosen under \ui{Apply as}.

\ui{Matches your setup} appears next to a preset when the enabled addons are
exactly the ones it names. This check always works as Replace: if one more
addon is enabled that the preset does not name, the preset no longer matches,
because Replace would disable that addon. A preset you have just saved matches
straight away, even though you never applied it.

Because the count follows the chosen mode and the mark does not, a preset can
match your setup and still show changes. In Disable mode, for example, it would
disable every addon it names.

\subsection{The three apply modes}

\begin{description}
  \item[Replace] Enables the preset's addons and disables everything else.
  \item[Accumulate] Enables the preset's addons and leaves the rest as it is.
  \item[Disable] Disables the preset's addons. Entries that the preset marks as
    disabled are ignored in this mode.
\end{description}

Before you apply, the \ui{Plan} panel splits the result into five groups:

\begin{itemize}
  \item \ui{To enable}: addons in the preset that are disabled now.
  \item \ui{To disable}: addons the preset disables, or that Replace disables
    because the preset does not name them.
  \item \ui{Already as the preset asks}: addons that stay as they are.
  \item \ui{Not found in the library}: addons the preset names that are no
    longer in the library. The rest of the preset is still applied.
  \item \ui{Requested, but not applied}: startup entries the preset asks for
    that could not be changed.
\end{itemize}

\shot{03b-presets-plan.png}{The plan of the ``Voo curto'' preset in Replace
mode. The sentence under the groups says the 4 addons outside the preset are
part of the 4 being disabled.}

\subsection{What Replace disables by omission}

In Replace mode, addons that are enabled now but not named in the preset are
disabled. They are listed under \ui{Disabled by Replace}, and \ui{Show them}
opens the list.

They are already part of the plan's count. If the plan says it disables 12 and
the list has 5 names, 12 addons are disabled in total: those 5, plus 7 the
preset itself disables.

\subsection{The return preset}

Just before applying a preset, the program saves the addons enabled at that
moment as the profile's return preset. \ui{Back to the previous set} applies it.
If the return preset cannot be saved, for example because the presets folder is
read-only, nothing is applied.

This is different from \ui{Undo last change}:

\begin{description}
  \item[Undo last change] Undoes what you just applied. It stops being
    available once you import an addon, work in the quarantine or change a
    category.
  \item[Back to the previous set] Applies the return preset like any other
    preset. It stays available after imports and quarantine work.
\end{description}

\subsection{Presets and startup entries}
\label{sec:presets-startup}

On the preset's \ui{Startup} tab, a checkbox lets the preset also control the
programs the simulator launches with itself. Checking it saves the startup
entries that are enabled at that moment. You can then enable or disable each
one there.

If startup management is off in \ui{Options}, the preset's startup entries are
not applied, and the program tells you how many were skipped. The addons are
applied as usual.

If an entry the preset names is no longer in the simulator's file, the program
lists it. It never adds startup entries.

\section{First setup}

The first time you open the program, it looks for simulator installations and
asks for two things:

\begin{description}
  \item[The simulator] Pick the installation for this profile from the list, or
    choose a folder yourself. If the folder does not look like the simulator's
    \where|Community| folder, the program says so and uses it anyway.
  \item[The library] Choose the folder outside the simulator where you keep
    your addons. Its subfolders become categories.
\end{description}

When you add a library with a long path, the program shows how many characters
the path already uses and compares it with a short path. Windows refuses some
operations, including the Recycle Bin, on paths longer than 260 characters.
The program does not block you: it warns when a path gets in the way, and
\ui{Diagnostics} shows the longest path in each library.

If you used MSFS Addons Linker before, open \ui{Options} and use
\ui{Import from MSFS Addons Linker}. It adds the libraries, categories and
presets that are not here yet. No files are moved or deleted.

\section{Library}

The \ui{Library} tab shows your addons as a tree, grouped by category. The
checkbox on each row enables or disables the addon.

\shot{01-library.png}{The library, with the categories open and the selected
addon's panel on the right.}

\subsection{The columns}

\begin{description}
  \item[Addon] The folder name.
  \item[Version] The version declared in the addon's \where|manifest.json|.
  \item[Enabled] Whether the addon has a link in a destination right now.
  \item[Destination] Where the link is, and whether the destination is pinned.
  \item[Link] Whether the link works. If its folder is missing, the link is
    broken.
\end{description}

A row can carry one of three marks: \ui{Broken link}, \ui{In conflict} or
\ui{Two copies}. The sections on destinations and quarantine explain them.

\subsection{Enabling and disabling in bulk}

Select several rows and use \ui{Enable selected} or \ui{Disable selected}. For
more than 10 addons, the program asks first. \ui{Undo last change} undoes the
whole batch.

If another addon of yours already uses the folder name in the destination, the
program offers to swap them: it disables that one and enables yours in a single
step, which you can undo in one go.

\subsection{Categories}

Right-click the tree to create, rename or delete a category, or to move addons
between categories. Only an empty category can be deleted.

\ui{Suggest categories} looks at folder names and \where|manifest.json| files
and proposes where to move each addon. The reliable rules come checked. The
livery rule is often wrong, so it comes unchecked: review it before applying.

The program puts a hidden \where|.fsorg-category| file in each category folder
it knows about, so an empty category is still recognized as one. It marks them
when you add the library, and \ui{Categories} in the options can mark them
again or remove every marker.

A folder with no addons inside and no \where|manifest.json| of its own counts as
an addon, so you can enable, move and delete it like any other. Converted
packages that the simulator writes, and scenery object libraries shipped
without a manifest, look like this.

\subsection{Pinning the destination}

Addons are linked in the profile's default destination. \ui{Always use}
followed by a destination pins an addon, or a whole category, to that
destination. If a pinned destination is removed from the profile, the program
tells you and uses the default destination until you choose another. The pin
itself is kept.

\subsection{Repointing}

If a link points to the wrong place, or to a folder that no longer exists,
\ui{Repoint to the library} fixes it. Only the link changes; the addon's files
are not touched.

\section{Destinations}

The \ui{Destinations} tab lists everything in the folders the simulator reads,
including folders you did not add through FS Organizer. Each entry has a
classification:

\shot{02-community.png}{The destinations, with the classification filter at the
top and the count of each kind at the bottom.}

\begin{description}
  \item[Managed] The simulator loads it from your library. This is the normal
    state.
  \item[External] A regular folder installed by another program, not in a
    library yet.
  \item[Divergent] The program that installed the addon put its own copy back,
    so there are two copies.
  \item[Vanished] The library copy is gone, removed by the other program. There
    is nothing to repair: the files are no longer on this computer.
  \item[Broken] The link points to a folder that does not exist.
  \item[Unavailable] The drive is not connected. It may come back, so the
    program does not offer to clean it up.
  \item[Unmanaged] A regular folder, not in a library yet.
  \item[Duplicated] Linked in more than one destination.
  \item[Substituted] Something wrote a regular folder over a link this program
    made. The simulator loads that folder, and enabling, presets and the culprit
    search no longer reach your library copy.
\end{description}

The filter at the top shows one classification at a time, and \ui{Select all
shown} selects exactly what is on screen.

\subsection{Importing another program's folder}

\ui{Import this folder} moves the folder into your library and leaves a link in
its place, so the addon becomes one of yours.

The other program will not know about the link. Its next update may write into
the link, or replace it with a regular folder and leave you with two copies.
When that happens, the entry shows as \ui{Divergent} and the program asks which
copy to keep.

To give the addon back to the program that installed it, use \ui{Delete from
the library}.

\section{Importing}

Importing copies a regular folder into the library and leaves a link where it
was. The original is removed only after the copy has been verified. The journal
records each of the five steps:

\shot{27-community-import.png}{The import dialog: how much will be copied, into
which library and category, and where the addon will end up.}

\begin{enumerate}
  \item Copy to a temporary folder.
  \item Verify the copy.
  \item Move the copy into place.
  \item Remove the original folder.
  \item Create the link in the destination.
\end{enumerate}

While the copy runs, a window shows the current folder, the progress and a
\ui{Cancel} button. Cancelling deletes nothing: what was already copied is kept,
so the import can be resumed later.

\subsection{Verify after copying}

On the \ui{Links} tab in \ui{Options}, you choose how the copy is verified.

\begin{description}
  \item[By structure] Compares the number and size of the files. This is the
    default and it is fast.
  \item[By hash] Compares the content of every file. It catches more, but
    imports become several times slower, especially with many small files.
\end{description}

If the copy does not match, the original is not removed.

\subsection{Unfinished imports}

If an import is interrupted, the program offers three choices the next time it
opens: resume the import, discard the partial copy, or leave it as it is. The
original folder is untouched in every case.

If the interrupted operation was a conflict resolution, the partial copy can
only be discarded, because both copies are still in place.

\section{Quarantine}

When you resolve a conflict between two copies of an addon, the copy you do not
keep goes to the quarantine instead of being deleted. Items stay there until
you restore or discard them.

For each item, the program records where it came from, in a file next to the
item. That record does not depend on the journal, so you can restore an item
even if the journal is lost. The table shows both sources, \ui{The record says}
and \ui{The Journal says}, and flags any disagreement.

\shot{05-quarantine.png}{The quarantine, with each item's origin and a flag
where the record and the journal disagree.}

\ui{Restore} puts each folder back where it came from. If something is already
there, the dialog shows both versions and you decide; replacing moves the
current folder to the quarantine. \ui{Discard} deletes the item permanently,
after asking.

\section{Deleting an addon from the library}

\ui{Delete} offers three ways to remove an addon, and tells you in advance when
one of them will not work.

\begin{description}
  \item[Move to the Recycle Bin] You can restore the addon from the Recycle Bin.
    This option is refused when the addon does not fit in that drive's Recycle
    Bin, or when it has a path longer than the 260 characters the Recycle Bin
    accepts. Windows may remove older items from the Recycle Bin to make room.
  \item[Delete permanently] This cannot be undone.
  \item[Give it back to the program that installed it] Only offered for an
    addon imported from another program's folder. The files go back to that
    folder and leave the library.
\end{description}

If the addon is enabled, it is disabled first.

\section{Simulator}

The \ui{Simulator} tab has two panels, and both change files that belong to the
simulator. Before changing one, the program saves a backup copy next to it.
Neither panel changes anything while the simulator is running. Both start
unmanaged, and until then the program does not read or write the file: click
\ui{Manage startup entries} or \ui{Manage the package list} to turn each one on.

\subsection{Startup entries}

The simulator's \where|EXE.xml| file, next to its \where|UserCfg.opt|, lists the
programs it launches with itself. The \ui{Startup entries} panel shows them and
lets you enable or disable each one without editing XML. The program only
enables or disables entries that already exist. It never adds, removes or
reorders them.

\shot{19-simulator-startup.png}{The startup entries, with the program each one
launches and its switch.}

When you disable an addon that has a startup entry, the program warns you:
otherwise the simulator would keep trying to launch a program that is no longer
there. You can disable both in one step, which you can undo in one go.

\subsection{Package list}

The simulator ships its own airports. When one of your airports covers the same
place as one of them, the \ui{Package list} panel shows the pair and lets you
disable the simulator's airport.

This changes one value in the simulator's package list. Nothing is added,
removed or reordered. An entry that does not have that value is left alone,
because an invalid file makes the simulator reset the whole list.

When two of your own addons cover the same airport, the program only shows the
pair. To keep just one, disable the other on the \ui{Library} tab.

\section{Documents}

The \ui{Documents} tab finds the PDFs that come with your addons and opens them
in the program. \ui{Scan the library} looks for PDFs in every addon. The result
is saved, so later scans are quick.

PDFs are split into \ui{Documents} and \ui{Charts}. Charts are grouped by
airport, and also by type when the addon includes a catalogue of them. One
addon can bring a manual and a hundred and fifty charts, so each group shows
its count first.

The reader has an outline, search, fit to width and bookmarks. A bookmark you do
not name is named after its page. \ui{Open in a separate window} moves the reader
out of the tab.

\shot{26-documents.png}{The documents tab: charts grouped by airport on the left,
the reader in the middle and the document outline on the right.}

On a chart, the mouse wheel zooms and dragging moves the page. On a document,
both start off. The two buttons on the toolbar switch them.

\section{Diagnostics}

The \ui{Diagnostics} tab shows what the program found on disk. Each section
shows when it was measured, and you can measure it again at any time.

\shot{09-diagnostics-size.png}{The size on disk section, with the total per
category and per addon, and when it was measured.}

\begin{description}
  \item[Destination entries] How many entries there are of each kind, and the
    longest path in each library.
  \item[Broken or unavailable] Broken links and entries on a disconnected
    drive, with a button to repair the broken links.
  \item[Quarantine] How much space the quarantine uses.
  \item[Size on disk] How much each category and addon takes up. It may take a
    few seconds; \ui{Stop} stops after the current addon.
  \item[Airports in the scenery] Scans each addon's scenery and groups addons
    by what it finds: an airport code, an airport record that cannot be read,
    navigation data instead of scenery, or no airport record.
  \item[Modules loaded] Reads the report the simulator writes after a slow
    load. The simulator does not report loading time per package, so this
    shows which module each package loaded and how much memory it uses.
\end{description}

\subsection{Finding the culprit}

When the simulator crashes and you do not know which addon is responsible,
\ui{Find the culprit} narrows it down by halves. In each round, the program
enables part of your addons, you launch the simulator and answer \ui{It crashed}
or \ui{It worked}, and half of the remaining suspects are ruled out.

\shot{31-diagnostics-bisection.png}{The culprit search, with the current round,
what has been ruled out, and the two answers.}

Before you start:

\begin{itemize}
  \item The search works on units. A unit is one addon, or a group of addons
    that must stay together, for example because they share a model folder. The
    panel says why a group is together, and offers to split it when that is
    safe.
  \item The first round disables all your addons, to find out whether the cause
    is among them at all.
  \item Your setup is saved before the first round and restored when the search
    ends, however it ends. If the program closes during the search, it offers
    to restore your setup the next time it opens.
  \item The search assumes a single culprit. If the crash only happens when two
    addons are enabled together, the result may point to the wrong one.
\end{itemize}

\section{Journal}

Every operation that changes files on disk adds a line to the \ui{Journal}.
Lines are only ever added, and the program never deletes the journal.

Each line shows when it happened, the operation, the addon, the library, the
source, the destination and the result. \ui{Failures only} shows just what went
wrong, and the search box looks in addon names, paths and operations.

\ui{Undo last change}, on the \ui{Library} tab, undoes the whole last change,
not a single line of it.

\section{Options}

The gear button in the header opens the options. They have five tabs.

\shot{10-options-profiles.png}{The profiles and libraries tab.}

\begin{description}
  \item[Profiles and libraries] Choose the profile in use, add or remove
    profiles and libraries, and change destinations. Removing a library deletes
    no files: its links keep working in the simulator, but FS Organizer stops
    managing them. \ui{Categories} shows how many category folders are marked
    and lets you mark them or remove the markers.
  \item[Links] The link type, and how imports are verified.
  \item[Updates] Automatic, notify, or manual. Automatic downloads new versions
    and installs them when you close the program.
  \item[Language] English or Brazilian Portuguese. The change is immediate.
  \item[About] The version, the licence and the credits.
\end{description}

\section{Where things are kept}

The program keeps its data in \where|%LOCALAPPDATA%\fs-organizer|.
\where|settings.json| holds the profiles, libraries, destinations and your
choices. \where|journal\operations.jsonl| is the journal. \where|presets| holds
one file per preset, plus each profile's return preset. \where|bisection| holds
the state of a culprit search in progress, and is deleted when the search ends.

None of your addons are stored there. If you delete that folder, you lose the
configuration and the presets, but no addon: the program finds them again by
reading the destinations and libraries.

The only thing the program writes outside that folder is the hidden
\where|.fsorg-category| marker in your category folders. \ui{Categories} in
the options removes all of them.

\end{document}
